\documentclass[a4paper]{article}
% Español (México) con XeLaTeX: fontspec para fuentes Unicode, polyglossia
% para la división silábica y los nombres del documento en la variante mexicana.
\usepackage{fontspec}
\usepackage{polyglossia}
\setdefaultlanguage[variant=mexican]{spanish}
% Los nombres chinos van como «pinyin (original)»: xeCJK da a los caracteres
% Han su propia fuente; sin ella XeLaTeX los omite con un aviso.
% FandolSong no cubre algunos caracteres de nombres (U+73FA); WenQuanYi Zen Hei sí.
\usepackage[AutoFallBack=true]{xeCJK}
\setCJKmainfont{FandolSong-Regular.otf}
\setCJKfallbackfamilyfont{\CJKrmdefault}{WenQuanYi Zen Hei}
% polyglossia no traduce los nombres de listings.
\AtBeginDocument{\providecommand{\lstlistingname}{}\renewcommand{\lstlistingname}{Listado}%
  \providecommand{\lstlistlistingname}{}\renewcommand{\lstlistlistingname}{Índice de listados}}
\usepackage{amsmath,amssymb}
\usepackage{graphicx}
\usepackage[margin=2.35cm]{geometry}
\usepackage[most]{tcolorbox}
\usepackage{etoolbox}
\usepackage{listings}
\usepackage{xcolor}
\usepackage{hyperref}
\usepackage{booktabs,longtable,tabularx,array}
\usepackage{subcaption}
\usepackage{float}
\usepackage{enumitem}
\usepackage{microtype}
\usepackage{fancyhdr}
\usepackage{needspace}

\definecolor{kbBack}{HTML}{F2F6FA}
\definecolor{kbFrame}{HTML}{9DB4CC}
\definecolor{kbTitle}{HTML}{52708F}
\definecolor{ibBack}{HTML}{FBF5E7}
\definecolor{ibFrame}{HTML}{C9A961}
\definecolor{ibTitle}{HTML}{7D5F1E}
\definecolor{wbBack}{HTML}{FCF2F0}
\definecolor{wbFrame}{HTML}{D6A096}
\definecolor{wbTitle}{HTML}{9E4F43}
\definecolor{dbBack}{HTML}{F4F7F3}
\definecolor{dbFrame}{HTML}{AEC2AC}
\definecolor{dbTitle}{HTML}{5F7F64}

\tcbset{softbox/.style={
  enhanced,breakable,boxrule=0.8pt,arc=2pt,left=3mm,right=3mm,
  top=2mm,bottom=2mm,coltitle=white,fonttitle=\bfseries,
  toptitle=0.6mm,bottomtitle=0.6mm
}}
\newtcolorbox{knowledgebox}[1]{softbox,colback=kbBack,colframe=kbFrame,colbacktitle=kbTitle,title={#1}}
\newtcolorbox{importantbox}[1]{softbox,colback=ibBack,colframe=ibFrame,colbacktitle=ibTitle,title={#1}}
\newtcolorbox{warningbox}[1]{softbox,colback=wbBack,colframe=wbFrame,colbacktitle=wbTitle,title={#1}}
\newtcolorbox{dialoguebox}[1]{softbox,colback=dbBack,colframe=dbFrame,colbacktitle=dbTitle,title={#1}}

\lstset{
  language=Python,basicstyle=\ttfamily\small,
  keywordstyle=\color{kbTitle}\bfseries,stringstyle=\color{wbTitle},
  commentstyle=\color{dbTitle},breaklines=true,frame=single,numbers=left,
  numberstyle=\tiny\color{gray},captionpos=b,columns=fullflexible,
  keepspaces=true,showstringspaces=false,extendedchars=true
}
\BeforeBeginEnvironment{lstlisting}{\Needspace{12\baselineskip}}

\hypersetup{colorlinks=true,linkcolor=kbTitle,urlcolor=kbTitle,citecolor=wbTitle}
\setlist{itemsep=0.25em,topsep=0.4em}
\setlength{\parindent}{2em}
\setlength{\parskip}{0.25em}
\newcolumntype{L}[1]{>{\raggedright\arraybackslash}p{#1}}

\providecommand{\notetitle}{Notas del curso: [tema del curso]}
\providecommand{\noteauthors}{Elaboradas a partir de los materiales públicos de Stanford CS25}
\providecommand{\notedate}{Versión reescrita en 2026}
\providecommand{\videochannel}{Stanford Online}
\providecommand{\videopublishdate}{2022}
\providecommand{\videoduration}{[duración del video]}
\providecommand{\videourl}{}
\providecommand{\videocoverpath}{cover.jpg}
\providecommand{\coursesite}{https://web.stanford.edu/class/cs25/}
\providecommand{\repourl}{https://github.com/hqhq1025/ai-course-notes}
\providecommand{\skillversion}{wdkns/wdkns-skills@39f1a04}

\newcommand{\figsource}[1]{\par\vspace{-0.3em}{\footnotesize\color{gray}\emph{Fuente: #1}}\par}
\newcommand{\teachervoice}[1]{\begin{knowledgebox}{Nota de clase}#1\end{knowledgebox}}
\newcommand{\term}[1]{\textbf{#1}}

\pagestyle{fancy}
\fancyhf{}
\fancyfoot[C]{\thepage}
\fancyfoot[R]{\footnotesize\color{gray}\href{\repourl}{AI Course Notes}}
\renewcommand{\headrulewidth}{0pt}
\renewcommand{\footrulewidth}{0pt}

\newcommand{\makecscover}{
\begin{titlepage}
\centering
\vspace*{0.7cm}
{\Large Stanford CS25 · Transformers United\par}
\vspace{0.8cm}
{\huge\bfseries \notetitle\par}
\vspace{0.6cm}
{\large \noteauthors\par}
\vspace{0.25cm}
{\large \notedate\par}
\vspace{0.9cm}
\ifdefempty{\videocoverpath}{}{
  \includegraphics[width=0.86\textwidth,height=0.43\textheight,keepaspectratio]{\videocoverpath}\par
}
\vfill
\begin{tcolorbox}[width=0.92\textwidth,colback=black!2!white,colframe=black!45,boxrule=0.7pt,arc=2pt]
\textbf{Autor o canal del video}: \videochannel\par
\textbf{Fecha de publicación}: \videopublishdate\par
\textbf{Duración del video}: \videoduration\par
\textbf{Sitio del curso}: \href{\coursesite}{\nolinkurl{\coursesite}}\par
\ifdefempty{\videourl}{}{\textbf{Enlace del video}: \href{\videourl}{\nolinkurl{\videourl}}\par}
\textbf{Normas de elaboración}: \skillversion\ + las normas source-first / teacher-voice / visual-QA de este repositorio
\end{tcolorbox}
\end{titlepage}
\tableofcontents
\newpage
}


\renewcommand{\notetitle}{Lecture 01: Transformers United — Attention, arquitectura y familias de modelos}
\renewcommand{\noteauthors}{Elaborado a partir de la introducción de Advay Pal, Chetanya Rastogi y Div Garg, y de cinco artículos originales}
\renewcommand{\notedate}{Curso de Fall 2021 · versión reescrita source-first de 2026}
\renewcommand{\videochannel}{Stanford Online}
\renewcommand{\videopublishdate}{2022-07-08}
\renewcommand{\videoduration}{22:43}
\renewcommand{\videourl}{https://www.youtube.com/watch?v=P127jhj-8-Y}
\renewcommand{\videocoverpath}{cover.jpg}

\newcommand{\lecturefigure}[3]{%
\begin{figure}[H]
  \centering
  \includegraphics[width=0.96\textwidth]{slides-images/#1}
  \caption{#2}
  \figsource{#3}
\end{figure}
}

\begin{document}
\makecscover

\section{Auditoría de fuentes: esto es un mapa del curso, no una cronología de modelos}

Esta sesión es la introducción de CS25 Transformers United V1, impartida por Advay Pal, Chetanya Rastogi y Div Garg. El curso se impartió en el otoño de 2021 y el video público se subió en julio de 2022. El video original dura solo 22 minutos y 43 segundos, pero cumple tres funciones: explicar cómo se organiza la speaker series, ofrecer la línea histórica mínima que va de attention a Transformer y, con GPT-3 y BERT, establecer las coordenadas de las familias decoder-only/encoder-only.

En el repositorio no se encontró un PDF independiente de las slides, pero el video muestra con claridad el deck original. Esta nota recupera del video oficial 17 slides con información didáctica; se omitieron a propósito las páginas que solo separan capítulos y los build states repetidos de la animación de self-attention. Las fórmulas ampliadas, los ejemplos calculados a mano y el pseudocódigo sirven para que la nota pueda estudiarse por sí sola; no significan que cada línea haya aparecido textualmente en clase.

\begin{warningbox}{Instantánea histórica y conocimiento actual}
«Where we are» y «The Future» reflejan la perspectiva de la clase de 2021. Esta nota las conserva como field snapshot y no presenta los juicios de entonces sobre el futuro como hechos actuales de 2026. Las ampliaciones técnicas posteriores solo sirven para explicar mecanismos duraderos y se mantienen separadas de lo que se dijo en clase.
\end{warningbox}

\subsection{La interfaz de aprendizaje de la speaker series}

CS25 V1 no sigue un libro de texto capítulo por capítulo, sino que invita a investigadores de distintos campos a explicar cómo el Transformer entra en sus propios problemas. Por eso la introducción funciona como un API contract: primero establece la terminología y la arquitectura comunes, y después cada ponente sustituye por su cuenta los datos, el objetivo, la mask, la tokenization y la evaluation. El lector debe preguntarse en todo momento «qué estructuras pertenecen al Transformer y cuáles son diseño propio del campo».

\lecturefigure{slide-01-course-page.jpg}{CS25 V1 organiza el conocimiento sobre Transformer mediante guest lectures de distintos campos.}{Video oficial 00:30; página del curso Stanford CS25 V1.}

\begin{knowledgebox}{Lectura de la figura: la página del curso no es una captura administrativa, sino una topología del conocimiento}
La página reúne a los instructors, al faculty advisor y el speaker schedule, lo que indica que la evidencia principal del curso proviene de la práctica de investigación de primera línea. Las sesiones siguientes pueden abarcar NLP, vision, reinforcement learning, generative modeling o biology, pero el lenguaje común sigue siendo la representación de tokens, los attention patterns, los objetivos de entrenamiento y el costo del sistema. Al estudiar el curso completo, conviene registrar primero qué interfaz cambió cada ponente y después comparar qué mecanismos se reutilizan entre campos.
\end{knowledgebox}

\lecturefigure{slide-02-title.jpg}{El objeto central de Transformers United son los modelos de aprendizaje profundo, no los robots que se transforman.}{Video oficial 00:55.}

\begin{knowledgebox}{Lectura de la figura: la delimitación del alcance detrás de la broma}
El título usa la imagen de los juguetes Transformers para que se recuerde, pero los ponentes de inmediato acotan el alcance al deep learning. La verdadera «transformación» es la transferibilidad del grafo de cómputo: el mismo núcleo de token-mixing puede procesar palabras, patches de imagen, acciones, frames de audio o secuencias biológicas; cuando la semántica de la entrada cambia, también cambian la representación y el objetivo. El curso no reduce todos los campos a una misma tarea, sino que estudia cómo un esqueleto unificado se combina con los sesgos inductivos de cada campo.
\end{knowledgebox}

\teachervoice{La apertura aclara a propósito que «no son robots que se transforman en autos» y luego destaca que el Transformer se extendió de NLP a CV y RL. Esta transición ligera conserva la verdadera thesis del curso: que el nombre del modelo sea el mismo no significa que la tarea sea la misma, pero un interaction mechanism reutilizable está conectando comunidades de investigación que antes estaban dispersas.}

\subsection{Quién imparte el curso y por qué abarca varios campos}

La formación de los tres instructores abarca production software, generative modeling, reinforcement learning, robotics y NLP. Esta combinación influyó en la selección de temas: el curso se ocupa tanto de las fórmulas como de la manera en que los modelos entran en sistemas reales; trata el lenguaje, pero también invita a ver la architecture como una herramienta de investigación transferible.

\lecturefigure{slide-03-instructors.jpg}{La formación de los tres instructores abarca ingeniería de software, NLP, modelado generativo, RL y robotics.}{Video oficial 01:30.}

\begin{knowledgebox}{Lectura de la figura: las diferencias de formación determinan las diferencias de preguntas}
La perspectiva de ingeniería de software pregunta por el despliegue, las interfaces y la mantenibilidad; la perspectiva de NLP se interesa en la representación del lenguaje, el preentrenamiento y la transferencia; la perspectiva de RL/robotics se interesa en la toma de decisiones secuenciales, la observabilidad parcial y la retroalimentación de las acciones. El Transformer no resuelve estos problemas automáticamente, pero ofrece un sequence interaction primitive unificado que permite a distintos campos compartir experiencia en optimización, escalamiento y aprendizaje de representaciones.
\end{knowledgebox}

\subsection{Tres objetivos de aprendizaje}

La introducción plantea explícitamente tres outcomes: entender cómo funciona el Transformer; entender por qué va más allá de NLP; descubrir nuevas líneas de investigación a partir de las guest talks. Corresponden, respectivamente, a mechanism, transfer y research taste. Memorizar solo $QK^\top V$ no basta para cumplir el segundo y el tercero; perseguir solo nombres de aplicaciones tampoco permite juzgar si la transferencia es razonable.

\lecturefigure{slide-04-learning-goals.jpg}{Los objetivos del curso son, en orden, la comprensión del mecanismo, la aplicación en distintos campos y nuevas líneas de investigación.}{Video oficial 02:05.}

\begin{knowledgebox}{Lectura de la figura: convertir los objetivos en una lista de verificación de tres niveles}
El nivel de mechanism exige explicar cómo interactúan los tokens, cómo la mask restringe la información y cómo la position entra en la representación; el nivel de application exige explicar los cambios en tokenization, objective, data y evaluation; el nivel de research exige identificar los modos de falla actuales y las hipótesis falsables. Cada sesión posterior puede repasarse con estos tres niveles, para evitar que las guest lectures se conviertan en un directorio de empresas o artículos.
\end{knowledgebox}

\subsection{Resumen de la sección}

El papel de Lecture 01 es establecer una interfaz común. El curso conecta varios campos mediante las guest talks; esta nota completa ese mapa con todas las diapositivas de la fuente, la voz del docente y derivaciones autocontenidas.
\section{Attention Timeline: de dónde viene el Transformer}

La historia del Transformer no puede contarse como «en 2017 se inventó de repente la self-attention». El Transformer hereda una larga acumulación de trabajo en sequence modeling, encoder-decoder, alignment y attention. La función didáctica más importante de la línea de tiempo es mostrar el Transformer como un rediseño de la ruta que sigue la información: de la compresión recursiva se pasa a la lectura según una consulta y, después, a la interacción directa dentro de la secuencia.

\subsection{De la recurrent memory a la attention}

Una \term{RNN} (recurrent neural network, red neuronal recurrente) actualiza su estado oculto de forma recursiva a lo largo del tiempo:
\[
h_t = f(x_t,h_{t-1}),
\]
donde $x_t$ es la $t$-ésima entrada, $h_t$ es el estado acumulado hasta la posición actual y $f$ suele implementarse con una unidad no lineal con parámetros. \term{LSTM} y \term{GRU} son variantes de RNN con compuertas, que usan gates explícitas para mitigar el desvanecimiento del gradiente y controlar qué memoria se conserva. Pueden codificar el orden, pero la información tiene que propagarse paso a paso a lo largo de una cadena larga, y el entrenamiento difícilmente puede paralelizarse por completo en la dimensión temporal.

\lecturefigure{slide-05-attention-timeline.jpg}{Attention, self-attention y los modelos preentrenados se ubican en una línea de tiempo de evolución continua.}{Video oficial 03:10; Vaswani et al. 2017.}

\begin{knowledgebox}{Lectura de la figura: la unidad de la línea de tiempo es la ruta de información}
Los primeros sequence models comprimían la historia en un recurrent state; en 2014--2015, la attention permitió que el decoder leyera con pesos varios encoder states; en 2017, el Transformer fue más allá y permitió que los tokens hicieran self-attend directamente entre sí; después, modelos preentrenados como BERT y GPT extendieron la misma arquitectura como interfaz general de representación o de generación. La tendencia principal no es que haya cada vez más nombres de modelos, sino que la señal de supervisión, el paralelismo del cómputo y la reutilización crecen al mismo tiempo.
\end{knowledgebox}

\teachervoice{El ponente llama a «Attention Is All You Need» el punto de inflexión de una idea sencilla de attention, pero antes muestra la prehistoria. Recordatorio de clase: una arquitectura exitosa suele recombinar mecanismos existentes en una interfaz más adecuada para el hardware, los datos y el objetivo de entrenamiento, en lugar de crear todos sus componentes desde cero.}

\subsection{El fixed-vector bottleneck de seq2seq}

Un modelo \term{seq2seq} (sequence-to-sequence) codifica la secuencia de entrada en una representación y luego el decoder genera la salida paso a paso. Los primeros enfoques solían comprimir toda la secuencia fuente en un solo \term{context vector} (vector de contexto) $c=h_n$. Cuando la entrada es larga, todos los significados de las palabras, el orden y las dependencias compiten por una capacidad limitada; además, el decoder debe recuperar de ese mismo $c$ la información que necesita en cada paso.

\lecturefigure{slide-06-prehistory.jpg}{RNN, seq2seq, LSTM y GRU pueden codificar memoria, pero enfrentan las secuencias largas y el context bottleneck.}{Video oficial 03:45.}

\begin{knowledgebox}{Lectura de la figura: no tomes «las RNN no funcionan» como una conclusión absoluta}
Las cruces rojas de la slide destacan la dificultad con contextos largos y con el acceso por contenido; no indican que las RNN fallen en todas las tareas. Las RNN con compuertas fueron, y siguen siendo, eficaces en modelos pequeños, en procesamiento en flujo y en tareas con un fuerte sesgo secuencial. La ventaja del Transformer proviene de rutas de información más cortas y del entrenamiento en paralelo, pero tiene un costo: la attention matrix, la memoria y un sesgo local más débil. Una comparación correcta debe especificar la longitud de la secuencia, la latencia, la escala de los datos y el hardware.
\end{knowledgebox}

\begin{importantbox}{La esencia del bottleneck}
Si en el paso $t$ el decoder solo puede leer el mismo $c$, no puede preguntar directamente «qué posición de la secuencia fuente es la más relevante para lo que se genera ahora». Lo esencial de la attention no es hacer más grande el vector, sino cambiar la interfaz a una query-conditioned retrieval: cada paso de generación puede producir sus propios pesos y leer posiciones distintas.
\end{importantbox}

\subsection{Resumen de la sección}

La prehistoria del Transformer es una ruta de acceso a la información: la RNN arrastra un estado de forma recursiva, seq2seq comprime la secuencia, la attention lee según lo que se necesita y la self-attention permite que cada token sea a la vez quien consulta y fuente de información.
\section{2021 Field Map: dónde estamos y qué falta}

Entender la arquitectura también exige entender por qué atrajo la atención de varios campos. Para 2021, el Transformer ya se había extendido de la generación y comprensión de lenguaje a vision, generative modeling, reinforcement learning y a sistemas relacionados con la biología estructural. Aquí no basta con enumerar nombres de aplicaciones; hay que preguntar cómo define cada campo el token, la position, la mask y el objective.

\subsection{La transferencia entre campos no es una copia simple}

En lenguaje, el token suele ser una subword; en imágenes, el token puede ser un patch; en RL se puede serializar la secuencia state/action/reward; las proteínas pueden modelarse con aminoácidos o con una residue representation. La misma fórmula de self-attention procesa relaciones completamente distintas en cada campo; por eso el preprocessing, el inductive bias y la evaluation determinan si la transferencia funciona.

\lecturefigure{slide-07-where-we-are.jpg}{En la clase de 2021 ya se ubicaba el Transformer en un mapa común de NLP, vision, RL y modelado generativo.}{Video oficial, 04:30.}

\begin{knowledgebox}{Lectura de la figura: primero identificar el token y después discutir la arquitectura}
En NLP, el token tiene un orden discreto de palabras y una gramática; un patch de vision además necesita posición bidimensional y estructura local; una trajectory de RL contiene la causalidad de las acciones y la recompensa; un generative image model puede trabajar sobre tokens de pixel, latent o codebook. Decir solo «todos usan Transformer» oculta las decisiones de modelado más importantes. Al leer las clases siguientes, conviene anotar primero cuatro columnas: token, position, mask, loss.
\end{knowledgebox}

\begin{warningbox}{AlphaFold2 no es «solo Transformer»}
La Transformer-like attention es un componente importante de AlphaFold2, pero el sistema también combina alineamientos de secuencias múltiples, pair representation, módulos geométricos e inferencia estructural. El éxito entre campos suele venir del diseño conjunto de la attention y la estructura propia del campo, no de aplicar sin cambios un modelo de texto a datos nuevos.
\end{warningbox}

\subsection{Los problemas persistentes en la Future slide}

Por un lado, la slide imagina aplicaciones como video understanding, finance y generative modeling; por otro, enumera los componentes que faltan: external memory, la reducción de la complejidad computacional y la alineación con los objetivos humanos. Estos problemas persisten porque el Transformer solo define cómo se mezcla la información; no aporta por sí mismo estado de largo plazo, cómputo de bajo costo ni objetivos correctos.

\lecturefigure{slide-08-future.jpg}{La expansión de aplicaciones y los componentes faltantes deben leerse a la vez.}{Video oficial, 05:55.}

\begin{knowledgebox}{Lectura de la figura: la columna derecha se parece más a una agenda de investigación que la izquierda}
Las aplicaciones de la izquierda responden «en qué tareas podría aportar valor»; los missing ingredients de la derecha responden «por qué la arquitectura por sí sola todavía no basta». La external memory se ocupa de la información persistente que excede el context; la efficiency, de la attention $O(n^2)$ y del costo de entrenamiento; la alignment, de los objetivos, la retroalimentación y las restricciones sociales. Toda proposal de aplicación debería corresponder al menos a una hipótesis de capacidad y a un modo de falla.
\end{knowledgebox}

\teachervoice{El ponente no describe el Transformer como un punto final, sino que coloca las aplicaciones junto a los componentes faltantes. Nota de clase: el criterio de investigación surge de ver al mismo tiempo «qué se puede hacer» y «por qué todavía no se hace bien», y no de dejar de preguntar después de una demo exitosa.}

\subsection{Resumen de la sección}

La unificación entre campos proviene de la interaction primitive, no de la homogeneización de las tareas. El mapa del futuro del curso deja la tokenización, la memory, la efficiency y la alignment como problemas pendientes, y ofrece así coordenadas para las guest lectures.
\section{Basic Attention: de la compresión fija a la lectura según la necesidad}

Antes de llegar a self-attention, conviene dominar la attention de codificador-decodificador. El codificador produce un conjunto de estados $h_1,\ldots,h_n$, uno por cada posición de origen. En el instante $t$, el decoder usa el estado $s_{t-1}$ para calcular la compatibilidad, los pesos normalizados y el context:
\[
e_{t,i}=v^\top\tanh(W_hh_i+W_ss_{t-1}),\qquad
\alpha_{t,i}=\frac{\exp(e_{t,i})}{\sum_j\exp(e_{t,j})},\qquad
c_t=\sum_i\alpha_{t,i}h_i.
\]
Aquí $e_{t,i}$ es el score entre la $t$-ésima query y la posición de origen $i$, $\alpha_{t,i}$ es el attention weight normalizado y $c_t$ es el resultado de la lectura en este paso. Esta additive attention al estilo de Bahdanau sustituye el context vector único por un context distinto en cada paso.

\subsection{Soft attention y hard attention}

La \term{soft attention} asigna pesos continuos a todas las posiciones candidatas y, por lo general, puede entrenarse de extremo a extremo mediante retropropagación. La \term{hard attention} selecciona o muestrea unas cuantas posiciones; el cálculo puede ser más disperso, pero la selección discreta normalmente no es diferenciable de forma directa y requiere policy-gradient, relaxation u otro estimator. Entre ambas no hay un orden fijo del tipo «la soft es más precisa y la hard es más rápida»: el resultado real depende de la tarea y del método de entrenamiento.

\lecturefigure{slide-09-part-soft-attention.jpg}{La attention visual temprana ya distinguía entre soft selection y hard selection.}{Video oficial 06:30; cita de Xu et al. 2015 en la diapositiva.}

\begin{knowledgebox}{Lectura de la figura: los pesos continuos y la selección discreta cambian la interfaz de entrenamiento}
La soft attention equivale a un promedio ponderado diferenciable sobre varias regiones, de modo que el gradiente llega a cada posición; la hard attention equivale a elegir una región o unas pocas, con una ruta de cálculo más definida, pero introduce varianza de muestreo y el problema de asignación de crédito. El ejemplo del gato en la figura no busca demostrar que el modelo «entiende al gato», sino mostrar que los pesos de atención pueden funcionar como variable intermedia para encaminar la información. Además, visualizar los pesos no equivale de manera automática a una explicación causal.
\end{knowledgebox}

\begin{warningbox}{El attention weight no tiene por qué coincidir con una explicación humana}
Un peso alto indica en qué medida el cálculo actual lee cierto value, pero el modelo también pasa por residual, multi-head, nonlinearity y las capas posteriores. Un solo heatmap no basta para afirmar cuáles son las razones completas del razonamiento del modelo; hace falta verificarlo con ablation, counterfactual o causal intervention.
\end{warningbox}

\subsection{Global attention y local attention}

La \term{global attention} asigna un score a todas las posiciones de origen en cada decoding step; la \term{local attention} solo asigna scores dentro de una ventana predicha o de un vecindario seleccionado. La global ofrece acceso completo, pero su costo crece con la longitud; la local introduce un sesgo hacia los vecinos cercanos y reduce el cálculo, aunque puede pasar por alto evidencia lejana. Las propuestas posteriores de sliding-window, block-sparse y routing attention pueden verse como extensiones de este espacio de diseño.

\lecturefigure{slide-10-local-global-attention.jpg}{La global attention lee toda la secuencia; la local attention solo lee una ventana local.}{Video oficial 07:10; Luong et al. 2015.}

\begin{knowledgebox}{Lectura de la figura: la ventana es un sesgo inductivo, no solo un recurso de optimización}
El global model supone que cualquier posición puede ser relevante; el local model supone que la mayoría de las dependencias se concentran en el vecindario. El tamaño de la ventana, la forma de predecir su centro y el manejo de los bordes cambian las relaciones que el modelo puede expresar. En documentos largos, audio o video, la estructura local suele existir de verdad, pero las referencias entre segmentos son igual de importantes; el sistema debe elegir entre rendimiento, memoria y las dependencias de la tarea, en lugar de tratar la sparse attention como un sustituto sin pérdidas.
\end{knowledgebox}

\begin{importantbox}{La primera reestructuración de la interfaz con attention}
La attention temprana todavía conservaba el recurrent encoder/decoder, pero cambió la forma de leer: de un fixed vector a una content-dependent weighted retrieval. El siguiente paso, la self-attention, consiste en que cada posición de una misma secuencia lea las demás posiciones mediante una interfaz similar.
\end{importantbox}

\subsection{Resumen de la sección}

La basic attention ya contiene los ejes de diseño centrales de los sistemas posteriores: continua o discreta, global o local, acceso completo o encaminamiento disperso. La innovación del Transformer no consiste en inventar los «pesos», sino en convertir la attention en la capa principal de cálculo sobre secuencias.
\section{Self-Attention: Query, Key, Value e interacción directa}

La \term{self-attention} (autoatención) hace que la secuencia de entrada produzca al mismo tiempo query, key y value. Cada posición compara su query con todas las key y luego combina los value según los pesos obtenidos. Su diferencia con la encoder-decoder attention no está en que la fórmula sea completamente distinta, sino en que query, key y value provienen de la misma secuencia, por lo que los token pueden establecer directamente dependencias internas.

\subsection{Scaled dot-product attention}

Sea la matriz de entrada $X\in\mathbb{R}^{n\times d_{model}}$; mediante proyecciones lineales se obtiene
\[
Q=XW_Q,\qquad K=XW_K,\qquad V=XW_V.
\]
La scaled dot-product attention es
\[
\operatorname{Attention}(Q,K,V)
=\operatorname{softmax}\!\left(\frac{QK^\top}{\sqrt{d_k}}+M\right)V.
\]
$n$ es el número de token, $d_k$ es la dimensión de key/query, $QK^\top$ da todos los score query-key, $\sqrt{d_k}$ evita que el dot product crezca demasiado cuando aumenta la dimensión, y $M$ es una \term{mask} (máscara) opcional que suma $-\infty$ a las posiciones a las que no se permite acceder, de modo que, después del softmax, su peso es prácticamente cero.

\lecturefigure{slide-11-self-attention-formulation.jpg}{La self-attention produce query, key y value a partir de la misma secuencia.}{Video oficial, 08:00; Vaswani et al. 2017.}

\begin{knowledgebox}{Lectura de la figura: un token cumple tres papeles a la vez}
La query expresa «qué necesito ahora», la key expresa «con qué índice puedo ser encontrado» y el value expresa «qué contenido aporto si me eligen». Las tres provienen de proyecciones lineales distintas, de modo que una misma palabra puede usar un subespacio para determinar la relevancia y otro para transmitir información. La salida sigue siendo un conjunto de token representation, no un único vector en el que se comprime toda la secuencia.
\end{knowledgebox}

\begin{importantbox}{Intuición con un cálculo a mano de tres token}
Supongamos que los puntajes escalados entre una query y tres key son $[2,1,0]$; después del softmax quedan aproximadamente en $[0.665,0.245,0.090]$. Si los values son $v_1,v_2,v_3$, la salida es $0.665v_1+0.245v_2+0.090v_3$. Esto no equivale a recuperar un registro de una base de datos, sino a una mezcla ponderada diferenciable; el entrenamiento aprende simultáneamente el score space y el value content.
\end{importantbox}

\begin{lstlisting}[caption={Pseudocódigo mínimo de la scaled dot-product attention}]
scores = (query @ key.transpose(-2, -1)) / sqrt(key_dim)
scores = scores + mask
weights = softmax(scores, dim=-1)
output = weights @ value
\end{lstlisting}

\subsection{Rutas de información y paralelismo}

En una RNN, para que el primer token influya en el último normalmente hay que pasar por $O(n)$ pasos recurrentes; en una sola capa de full self-attention, dos posiciones cualesquiera interactúan directamente a través de una attention edge. Durante el entrenamiento, todas las query pueden calcularse en paralelo mediante multiplicación de matrices, lo cual se ajusta al alto rendimiento de los accelerator. El costo es que la score matrix explícita tiene $n\times n$ entradas.

\lecturefigure{slide-12-self-attention-diagram.jpg}{La self-attention despliega las relaciones token-to-token como conexiones directas que pueden calcularse en paralelo.}{Video oficial, 10:30; animación de clase, final state.}

\begin{knowledgebox}{Lectura de la figura: observe las conexiones, no solo los recuadros}
En la figura, los token de entrada pasan por proyecciones lineales que producen varios grupos de representaciones, y cada posición de salida agrega información proveniente de las demás posiciones. En comparación con una recurrent chain, la ruta de dependencia es más corta; en comparación con la convolution, las conexiones no son un kernel local fijo, sino que las determina dinámicamente el contenido. Lo que la figura no muestra es la división de trabajo entre attention head, la residual path ni el apilamiento de capas; por lo tanto, explica el routing, pero no representa un Transformer block completo.
\end{knowledgebox}

\begin{warningbox}{Una conexión directa no implica aprender automáticamente las relaciones correctas}
La self-attention permite que cualquier token interactúe con cualquier otro, pero el modelo sigue necesitando datos, objective, position, optimization y suficientes capas para aprender patrones útiles. Una capacidad alta también puede ajustarse a shortcut, spurious correlation o filtración de información; la accesibilidad es una condición para la capacidad, no una garantía de corrección.
\end{warningbox}

\subsection{Multi-head attention}

La \term{multi-head attention} (atención multicabeza) ejecuta en paralelo varias attention de menor dimensión:
\[
\operatorname{head}_r=\operatorname{Attention}(QW_r^Q,KW_r^K,VW_r^V),
\qquad
\operatorname{MHA}=\operatorname{Concat}(\operatorname{head}_1,\ldots,\operatorname{head}_H)W^O.
\]
Distintas head pueden aprender patrones locales, de largo alcance, sintácticos o posicionales, pero no se garantiza que cada head tenga un significado claro para una persona. La función principal de las múltiples cabezas es proporcionar varios subespacios de representación y varios routing channel en paralelo.

\begin{knowledgebox}{Términos clave: Q/K/V, head y mask}
\begin{tabularx}{\textwidth}{L{0.18\textwidth} X}
\toprule
Término & Problema que resuelve y relación con el curso \\
\midrule
Query & Expresa la necesidad de lectura de la posición actual; determina a qué key interroga. \\
Key & Proporciona el índice de coincidencia de las posiciones candidatas; junto con la query produce el score. \\
Value & Contenido que se mezcla según los pesos; la relevancia y el contenido transmitido pueden usar subespacios distintos. \\
Head & Proyección y attention channel independientes; amplían los tipos de relación y la capacidad de representación. \\
Mask & Impone límites a la información; la padding mask ignora las posiciones vacías y la causal mask impide ver el futuro. \\
\bottomrule
\end{tabularx}
\end{knowledgebox}

\subsection{Resumen de la sección}

El núcleo de la self-attention es la coincidencia aprendible entre query y key y la mezcla de value. Acorta las rutas de dependencia y aumenta el paralelismo del entrenamiento, pero genera una score matrix cuadrática y requiere position, mask, residual y normalization para convertirse en una arquitectura estable.
\section{Componentes necesarios del Transformer Block}

Escribir solo la attention equation deja fuera la mayor parte del modelo real. El orden de los tokens, la optimización en profundidad, la transformación no lineal y los límites de la información quedan a cargo, respectivamente, de la positional representation, de residual/normalization, de la feed-forward network y de la mask. No son adornos: son la base que permite que attention sea entrenable, apilable y aplicable a distintas tareas.

\subsection{Position, residual, normalization y masking}

La self-attention pura tiene permutation equivariance respecto del orden de la entrada: si se reordenan los tokens a la vez, la salida se reordena de la misma manera, pero el modelo no conoce el significado absoluto de «el primero» y «el último». La \term{positional encoding} (codificación posicional) inyecta el orden en la token representation. El Transformer original usa seno y coseno:
\[
PE(pos,2i)=\sin\!\left(pos/10000^{2i/d}\right),\qquad
PE(pos,2i+1)=\cos\!\left(pos/10000^{2i/d}\right).
\]
$pos$ es la posición, $i$ es el índice del canal y $d$ es la dimensión del modelo. Los modelos modernos también usan con frecuencia position learned, relative o rotary, pero el objetivo común es expresar el orden relativo o absoluto.

\lecturefigure{slide-13-ingredients.jpg}{Position, normalization, residual connection y masking forman juntos un Transformer block utilizable.}{Video oficial, 12:30; Vaswani et al. 2017.}

\begin{knowledgebox}{Lectura de la figura: cada ingredient corrige un modo de falla distinto}
Position corrige la ausencia de orden; la residual connection permite que la capa aprenda incrementos y ofrece un atajo para el gradiente; la layer normalization controla la escala de cada token representation; la feed-forward network agrega, en cada posición de forma independiente, una transformación no lineal entre canales; la mask establece qué información es legítimamente visible. Si se elimina cualquiera de estos componentes, el modelo quizá siga funcionando, pero ya no tendrá la estabilidad de entrenamiento ni la semántica de tarea de la arquitectura original.
\end{knowledgebox}

La \term{residual connection} (conexión residual) suele escribirse como $y=x+F(x)$, de modo que el block aprende una corrección respecto de la entrada. La \term{layer normalization} (normalización por capa) estandariza la dimensión de features de cada token y facilita la optimización de redes profundas. La position-wise feed-forward network es
\[
\operatorname{FFN}(x)=W_2\sigma(W_1x+b_1)+b_2,
\]
que comparte parámetros entre todas las posiciones pero calcula cada una de forma independiente: attention se encarga del token mixing y la FFN del channel mixing.

\begin{warningbox}{La mask es un límite semántico, no solo un parámetro de rendimiento}
Si un causal language model omite la causal mask, durante el entrenamiento verá la información futura a la derecha del token objetivo; la loss puede verse anormalmente buena, pero la generación resultará completamente inútil. Si la padding mask es incorrecta, los símbolos de relleno se tratarán como contenido real. Cualquier implementación de attention debe probar primero la matriz de visibilidad y después el throughput.
\end{warningbox}

\subsection{Resumen de la sección}

El Transformer block combina token mixing y channel mixing, y mediante position, mask, residual y normalization conserva el orden, la causalidad y la capacidad de entrenamiento. Recordar esta división de funciones es más útil que memorizar un diagrama de arquitectura.
\section{Encoder, decoder y las tres familias de modelos}

El Transformer original se diseñó para la traducción automática y consta de un encoder y un decoder. Ambos apilan bloques de attention y FFN, pero difieren en el acceso a la información. El \term{encoder} (codificador) suele permitir que cada token de origen lea la secuencia de origen en ambas direcciones; el \term{decoder} (decodificador), al generar el $t$-ésimo token de destino, solo puede leer el prefijo de destino y, mediante cross-attention, lee las encoder outputs.

\subsection{Flujo de datos encoder-decoder}

Durante el entrenamiento, los source tokens pasan por el encoder y producen la memory; los target tokens, desplazados a la derecha, entran en el masked decoder. La self-attention del decoder procesa el prefijo ya generado; la \term{cross-attention} usa el decoder state como query y la encoder memory como key/value, y al final una softmax predice el siguiente token. \term{autoregressive} (autorregresivo) significa que la probabilidad conjunta se descompone como
\[
p(y\mid x)=\prod_{t=1}^{m}p(y_t\mid y_{<t},x).
\]

\lecturefigure{slide-14-encoder-decoder.jpg}{El Transformer original realiza seq2seq con la encoder memory, el masked decoder y la cross-attention.}{Video oficial 13:40; Vaswani et al. 2017.}

\begin{knowledgebox}{Lectura de la figura: no hay que mezclar los tres tipos de attention}
La self-attention del encoder lee todo el source; la masked self-attention del decoder solo lee el target prefix; la cross-attention permite que la target query lea la encoder memory. En la figura, el target desplazado a la derecha garantiza que el modelo no pueda ver directamente la respuesta actual. El entrenamiento puede calcular en paralelo todas las target positions, pero la máscara de cada position conserva la semántica autorregresiva; la inferencia sigue requiriendo avanzar token por token o recurrir a métodos especializados de decodificación en paralelo.
\end{knowledgebox}

\begin{importantbox}{Las tres familias principales}
\begin{enumerate}
  \item \textbf{Encoder-only}: representación bidireccional, adecuada para clasificación, recuperación, etiquetado y representation learning; ejemplo representativo: BERT.
  \item \textbf{Decoder-only}: causal next-token modeling, adecuada para la generación abierta y para una prompt interface unificada; ejemplo representativo: la serie GPT.
  \item \textbf{Encoder-decoder}: source y target separados, adecuada para traducción, resumen y generación condicionada; ejemplos representativos: el Transformer original y T5.
\end{enumerate}
Elegir una familia es, ante todo, elegir los límites de la información y la interfaz de la tarea, no solo la cantidad de parámetros.
\end{importantbox}

\subsection{Resumen de la sección}

El encoder, el decoder y la cross-attention definen de dónde proviene la información y cuándo es visible. La diferencia entre GPT y BERT puede deducirse directamente de este diagrama de la arquitectura original, en lugar de memorizarse como dos nombres de modelos aislados.
\section{Ventajas y costos: por qué el Transformer tiene éxito y por qué sigue siendo difícil}

El éxito del Transformer suele atribuirse a las dependencias de largo alcance y al entrenamiento en paralelo, pero ambas ventajas tienen condiciones. La full attention ofrece rutas de información cortas, pero genera $n^2$ score; el cálculo matricial se adapta bien a GPU/TPU, pero exige lotes grandes, memoria y conjuntos de datos de gran escala; un sesgo inductivo local más débil aumenta la generalidad, pero también puede incrementar la cantidad de muestras necesarias.

\subsection{Complejidad y límites del sistema}

En una sola capa de full self-attention, el cálculo de los score es de aproximadamente $O(n^2d)$ y la score memory, de aproximadamente $O(n^2)$; una sola capa de RNN suele costar alrededor de $O(nd^2)$ y requiere $O(n)$ sequential steps. Cuál de las dos es más rápida depende de $n$, $d$, el lote, el kernel, el memory bandwidth y el hardware. Una tabla de complejidad no equivale directamente a la wall-clock latency.

\lecturefigure{slide-15-good-bad.jpg}{El modelado de largo alcance y el paralelismo del Transformer vienen acompañados de una attention cuadrática, de la demanda de datos y de un sesgo local más débil.}{Video oficial, 16:10.}

\begin{knowledgebox}{Lectura de la figura: las ventajas y las desventajas provienen del mismo diseño}
La interacción directa entre cualquier par de token permite capturar dependencias de largo alcance, pero también produce una dense pairwise matrix; no depender de la recurrence permite entrenar en paralelo, pero también reduce el sesgo de orden incorporado; un block uniforme facilita el escalamiento, pero puede requerir más datos para aprender la estructura del dominio. Al leer la figura, no hay que tratar las dos columnas como listas independientes, sino buscar las dos caras de una misma decisión de diseño.
\end{knowledgebox}

\begin{warningbox}{FlashAttention no vuelve lineal la complejidad matemática}
La fused/tiled attention puede reducir las lecturas y escrituras en \term{HBM} (High Bandwidth Memory, memoria de video de alto ancho de banda de la GPU) y la materialización de matrices intermedias, lo que mejora notablemente la velocidad y el uso de memoria en la práctica; sin embargo, el cálculo pairwise de la full attention estándar sigue siendo cuadrático. La algorithmic sparsity, el low-rank, los state-space o la recurrence son enfoques distintos; la optimización del sistema y la complejidad algorítmica deben describirse por separado.
\end{warningbox}

\begin{knowledgebox}{Términos clave: inductive bias y sample efficiency}
El \term{inductive bias} (sesgo inductivo) es la estructura que el modelo prefiere antes de ver los datos; por ejemplo, el compartir pesos local e invariante a traslaciones de la convolution, la recursión temporal de la RNN o la direccionalidad de la causal mask. Un sesgo menor ofrece mayor generalidad, pero el modelo puede necesitar más datos para aprender regularidades que podrían haberse codificado a mano. La \term{sample efficiency} describe la cantidad de datos necesaria para alcanzar el desempeño objetivo, no el número de muestras procesadas por segundo.
\end{knowledgebox}

\subsection{Resumen de la sección}

Las ventajas y los costos del Transformer tienen el mismo origen: la conexión global dinámica, el paralelismo matricial y los supuestos débiles sobre el dominio amplían la escalabilidad, pero también traen un costo cuadrático, demanda de datos y complejidad del sistema. El juicio de ingeniería debe considerar al mismo tiempo el algoritmo y el hardware.
\section{GPT-3: Decoder-Only e In-Context Interface}

La parte de aplicaciones usa primero GPT-3 para mostrar la vía decoder-only. El modelo apila causal decoder blocks y se preentrena con next-token prediction sobre texto a gran escala. El objetivo de entrenamiento es simple, pero unifica muchas tareas como «dado un prefijo, predecir una continuación razonable». La instruction, los example y el context del prompt pasan a formar un mismo token stream.

\subsection{Autoregressive preentrenamiento y prompting}

Dada una secuencia de tokens $x_1,\ldots,x_T$, el modelo de lenguaje maximiza
\[
\sum_{t=1}^{T}\log p_\theta(x_t\mid x_{<t}).
\]
La causal mask garantiza que la posición $t$ solo vea el prefijo. \term{in-context learning} (aprendizaje en contexto) se refiere a que el modelo cambia su comportamiento en la tarea actual a partir de las instrucciones o los ejemplos dentro del prompt, sin actualizar sus parámetros. No es gradient-based training tradicional ni garantiza obtener reglas estables; el prompt format, el example order y el distribution shift influyen en el resultado.

\lecturefigure{slide-16-gpt3.jpg}{GPT-3 forma una interfaz de generación general con preentrenamiento causal decoder-only y prompt examples.}{Video oficial 18:30; Brown et al. 2020.}

\begin{knowledgebox}{Lectura de la figura: distinguir preentrenamiento, fine-tuning e in-context learning}
El preentrenamiento actualiza los parámetros para que el modelo aprenda una token distribution amplia; el fine-tuning sigue actualizando los parámetros con datos de la tarea; el in-context learning solo cambia el contexto de entrada y los parámetros permanecen iguales. Los ejemplos few-shot de la figura muestran que la task queda codificada en el prompt, no prueban que el modelo realmente ejecute un descenso de gradiente interno. Al evaluar, hay que comparar zero-shot, few-shot y el fine-tuned baseline, y controlar el token budget y la contamination.
\end{knowledgebox}

\teachervoice{El ponente usa GPT-3 para enfatizar «usar la red preentrenada como generator»: con una descripción de la tarea y unos cuantos example, produce una completion. La nota de clase es el cambio de interfaz: muchas tareas pasan de modificar el código del modelo a diseñar el contexto; esto aumenta la reutilización, pero también traslada los problemas de confiabilidad al prompt, a la evaluation y a los guardrails.}

\begin{warningbox}{El next-token objective no equivale a alineación con los hechos o con la intención}
Lo que optimiza el modelo de lenguaje es la likelihood condicional de los tokens, no «decir solo información verdadera» ni «cumplir el objetivo real del usuario». La escala y el prompting pueden mejorar el desempeño en las tareas, pero la exactitud factual, la calibration, la seguridad y la tool-use correctness requieren datos, entrenamiento y controles de sistema adicionales.
\end{warningbox}

\subsection{Resumen de la sección}

GPT-3 convierte el Transformer decoder-only en una prefix-to-continuation interface general. Muestra la escala del preentrenamiento y el in-context behavior, pero no elimina la definición de la tarea, la evaluación ni los límites de seguridad.
\section{BERT: Encoder-Only y Bidirectional Representation}

BERT muestra otra vía: usa solo bloques encoder, permite que cada token lea el contexto en ambas direcciones y aprende una representation mediante masked language modeling. Su interfaz principal no es la generación autorregresiva abierta; en cambio, ofrece contextual embeddings a los que se les puede aplicar fine-tuning para tareas posteriores como clasificación, span prediction, recuperación y token labeling.

\subsection{Masked language modeling}

El \term{masked language modeling} (MLM, modelado de lenguaje enmascarado) elige al azar algunas posiciones de la entrada y le pide al modelo que prediga el token original. Como las posiciones objetivo están enmascaradas, el encoder puede usar el contexto de ambos lados sin leer directamente la respuesta. El BERT original también usaba next sentence prediction (NSP). Los modelos posteriores resolvieron de distintas maneras si NSP era necesaria, pero lo que se ha mantenido como núcleo es la representation bidireccional del encoder.

\lecturefigure{slide-17-bert.jpg}{BERT se apoya en una representación bidireccional encoder-only y en MLM para las tareas posteriores de comprensión.}{Video oficial 20:30; Devlin et al. 2018.}

\begin{knowledgebox}{Lectura de la figura: la diferencia entre BERT y GPT es, ante todo, de visibilidad}
En BERT, cada token puede leer a la vez el contexto izquierdo y el derecho, por lo que es adecuado para aprender la contextual representation de una oración completa. En GPT, cada token causal solo puede leer el prefijo a su izquierda, lo que concuerda de forma natural con la generación paso a paso. El fine-tuning de BERT suele agregar una task head y actualizar los parámetros; con el prompting de GPT, la tarea puede describirse dentro de una interfaz de generación unificada. Ambos provienen del Transformer, pero difieren en la mask, el objective y la product interface.
\end{knowledgebox}

\begin{importantbox}{Comparación entre encoder-only y decoder-only}
\begin{tabular}{L{0.18\textwidth} L{0.36\textwidth} L{0.36\textwidth}}
\toprule
Dimensión & Encoder-only (BERT) & Decoder-only (GPT) \\
\midrule
Visibilidad & self-attention bidireccional & causal left-to-right attention \\
Objetivo de preentrenamiento & masked token prediction & next-token prediction \\
Salida típica & contextual representation & token distribution / generation \\
Interfaz de la tarea & fine-tuning, linear head, embedding & prompting, completion, fine-tuning \\
Ventaja & tareas intensivas en comprensión y representación & interfaz unificada de generación e interacción \\
\bottomrule
\end{tabular}
\end{importantbox}

\teachervoice{El ponente cierra la introducción con BERT y enfatiza el bloque encoder, la masked prediction y el fine-tuning para tareas posteriores. Este orden es importante: primero se separan dos direcciones a partir del encoder-decoder original y después se usan GPT y BERT para explicar las familias de arquitecturas, en lugar de deducir el mecanismo a partir de los nombres comerciales.}

\subsection{Resumen de la sección}

El núcleo de BERT es la representation bidireccional del encoder y la masked prediction. Comparte con GPT el bloque de self-attention, pero sirve a interfaces de tarea distintas mediante una mask y un objective diferentes.
\section{Cómo leer el resto de CS25: cuatro interfaces y una lista de verificación}

El valor final de la introducción no es detener la historia en GPT-3/BERT, sino ofrecer un lenguaje de comparación para el resto del curso. Ya sea que el ponente hable de visión, RL, modelos generativos, biology o modelos de lenguaje más grandes, todo puede desglosarse con cuatro preguntas: token interface, interaction pattern, training objective y deployment contract.

\subsection{Método de lectura de las cuatro interfaces}

\begin{importantbox}{Cada sesión responde cuatro preguntas}
\begin{enumerate}
  \item \textbf{Token interface}: ¿en qué unidades se divide la entrada y cómo se representan la position y la modality?
  \item \textbf{Interaction pattern}: ¿full, local, causal, cross, sparse o recurrent?
  \item \textbf{Training objective}: ¿next-token, mask, contrastive, control, ranking o multi-task?
  \item \textbf{Deployment contract}: ¿la salida se usa para generation, representation, decision, retrieval o action? ¿Cómo se evalúan las fallas?
\end{enumerate}
\end{importantbox}

\teachervoice{Al final, los tres instructores dirigen al público hacia la serie completa de guest-speakers. Nota de clase: la introducción solo aporta un vocabulary común; el aprendizaje real ocurre al ver cómo los ponentes siguientes modifican el token, la mask, el objective y el system boundary.}

\begin{knowledgebox}{Plantilla de pregunta de investigación}
Primero elige una falla del dominio: contexto largo, pocos datos, latencia en tiempo real, decisiones causales o restricciones de seguridad; después señala qué interfaz del Transformer estándar causa la limitación; propone un cambio verificable; presenta el baseline, la ablation y el presupuesto de recursos; por último, explica qué no pueden demostrar las figuras y tablas. Así, las «aplicaciones de Transformer» dejan de ser un simple cambio de conjunto de datos y se convierten en diseño de investigación.
\end{knowledgebox}

\subsection{Resumen de la sección}

El resto de CS25 no debe leerse como un catálogo de modelos. El método de las cuatro interfaces ayuda a distinguir el esqueleto general, la adaptación al dominio, la señal de entrenamiento y la responsabilidad del producto, y convierte cada guest talk en un diseño de sistema comparable.
\section{Síntesis y ampliación}

Esta sesión construye una visión mínima y completa del Transformer:

\begin{enumerate}
  \item \textbf{La historia no es una ruptura}: RNN/seq2seq dejan expuesto el fixed-vector bottleneck, la attention temprana ofrece lectura bajo demanda y la self-attention la convierte en la capa principal de secuencia;
  \item \textbf{La attention es routing}: la query expresa la necesidad, la key ofrece el índice de coincidencia y el value aporta el contenido que se transmite;
  \item \textbf{El espacio de diseño ya existía}: soft/hard y global/local controlan, respectivamente, la selección diferenciable y el alcance del acceso;
  \item \textbf{La self-attention acorta las rutas}: los token pueden interactuar directamente y entrenarse en paralelo, pero la full score matrix tiene un costo cuadrático;
  \item \textbf{Un block completo no es solo attention}: position, mask, residual, normalization y FFN resuelven, respectivamente, el orden, la visibilidad y la optimización en profundidad;
  \item \textbf{Las familias se definen por las fronteras de información}: encoder-only produce representaciones bidireccionales, decoder-only genera de forma causal y encoder-decoder procesa seq2seq condicional;
  \item \textbf{GPT-3 cambió la interfaz de las tareas}: el preentrenamiento decoder-only permite expresar las tareas mediante prompt y example;
  \item \textbf{BERT reforzó la interfaz de representación}: el MLM encoder-only aprende contextual representation bidireccional;
  \item \textbf{Pasar a otros dominios no es copiar}: cada dominio debe redefinir token, position, mask, objective y evaluation;
  \item \textbf{La arquitectura no resuelve por sí sola los problemas futuros}: memory, efficiency, alignment y deployment reliability siguen requiriendo un diseño propio.
\end{enumerate}

\begin{importantbox}{Ejercicio de cálculo a mano y de diseño}
Dados $Q,K,V$ para tres token, calcula a mano el score, el scaling, la softmax y el output; después agrega por separado una padding mask y una causal mask, y explica qué pesos deben valer cero. Luego elige uno de estos casos: imágenes, RL trajectory o secuencias de proteínas; define token, position, attention pattern, objective y metric, y explica por qué usarías encoder-only, decoder-only o encoder-decoder.
\end{importantbox}

\begin{warningbox}{Autoevaluación final}
Si solo puedes repetir «el Transformer usa attention», todavía no dominas la interfaz. Debes poder explicar en qué difieren los orígenes de la basic attention y de la self-attention, cómo la mask define las fronteras de información, por qué GPT y BERT comparten el block pero sirven a tareas distintas, y si $O(n^2)$ es una propiedad de la estructura matemática o un cuello de botella de una implementación concreta.
\end{warningbox}

\subsection{Lecturas adicionales}

{\small
\begin{itemize}[nosep,leftmargin=*]
  \item \href{https://arxiv.org/abs/1409.0473}{\textbf{Bahdanau et al. (2014)}}: additive attention en encoder-decoder y alignment.
  \item \href{https://arxiv.org/abs/1508.04025}{\textbf{Luong et al. (2015)}}: global/local attention y multiplicative scoring.
  \item \href{https://arxiv.org/abs/1706.03762}{\textbf{Vaswani et al. (2017)}}: scaled dot-product, multi-head y el Transformer encoder-decoder original.
  \item \href{https://arxiv.org/abs/2005.14165}{\textbf{Brown et al. (2020)}}: GPT-3 y la evaluación few-shot/in-context.
  \item \href{https://arxiv.org/abs/1810.04805}{\textbf{Devlin et al. (2018)}}: BERT, MLM, NSP y la interfaz de fine-tuning.
  \item \href{https://web.stanford.edu/class/cs25/past/cs25-v1/}{\textbf{Stanford CS25 V1}}: la speaker series original y el programa del curso.
\end{itemize}
}

\end{document}
